\uuid{IWbb}
\titre{Espaces vectoriels}
\auteur{hanani ; mzari}
\organisation{exo7}
\licence{CC BY-NC-SA 4.0}
\datecreate{2020-01-20}
\niveau{L1}
\module{Algèbre}
\chapitre{Espace vectoriel}
\sousChapitre{Définition, sous-espace}
\difficulte{2}
\theme{Espaces vectoriels}
\exo7qcmid{lille-2.espaces-vectoriels.008}

\contenu{
\begin{question}
\qclassification{106.01}
\qtags{motcle=Strucutre de sous-espace vectoriel de $\Rr^n$}
Soit $E=\{(x,y) \in \Rr^2 \,  ; \;  e^xe^y=1\}$, muni des opérations usuelles. Quelles sont les assertions vraies ?
\begin{answers}
\bad{$E=\{(x,y) \in \Rr^2 \,  ; \;  x=y \}$.}
\good{$E=\{(x,y) \in \Rr^2 \,  ; \;  x=-y \}$.}
\bad{$E$ est vide.}
\good{$E$ est un espace vectoriel.}
\end{answers}
\begin{explanations} $E=\{(x,y) \in \Rr^2 \, ; \;  x+y=0\}$ est un sous-espace vectoriel de $\Rr^2$.
\end{explanations}
\end{question}
}
